\section{Experimental protocol and preregistration}
\label{sec:protocol}

\subsection{Datasets and the development/held-out split}
Six public binary clinical datasets were used (Table~\ref{tab:datasets}). All come
from OpenML~\cite{vanschoren2014} with the file md5 published by OpenML, except
chronic kidney disease (UCI), whose sha256 was pinned at first download. Two
\emph{development} datasets, Pima diabetes and Cleveland heart
disease~\cite{detrano1989}, were used for every design decision, including the pilot
that led to the refinements in Section~\ref{sec:method}. Four \emph{held-out}
datasets (WDBC breast cancer, ILPD liver disease, thyroid ``sick'' and CKD) were
selected from metadata only. The data loader refused to read them until the
preregistration was locked. Positive classes were confirmed from the official dataset
documentation and checked on load against the registered sample size and number of
positives. Physiologically impossible zeros in Pima are a fixed row-wise rule and are
set to missing before splitting.

\begin{table}[t]
\centering
\scriptsize
\setlength{\tabcolsep}{3pt}
\caption{Datasets. Development datasets were used for all design decisions and the pilot; held-out datasets were unlocked only after the preregistration was locked. Features: numeric / binary / nominal before one-hot encoding.}
\label{tab:datasets}
\begin{adjustbox}{max width=\linewidth}
\begin{tabular}{llrrrl}
\toprule
Dataset & Role & $n$ & Prev. & Features & Source \\
\midrule
Pima & dev & 768 & 0.349 & 8/0/0 & OpenML 37 \\
Cleveland & dev & 303 & 0.455 & 6/3/4 & OpenML 49 \\
WDBC & held-out & 569 & 0.373 & 30/0/0 & OpenML 1510 \\
ILPD & held-out & 583 & 0.714 & 9/1/0 & OpenML 1480 \\
Thyroid & held-out & 3772 & 0.061 & 7/20/2 & OpenML 38 \\
CKD & held-out & 400 & 0.625 & 14/10/0 & UCI 336 \\
\bottomrule
\end{tabular}
\end{adjustbox}
\end{table}


\subsection{Preregistration}
Before any held-out file was downloaded, a preregistration was committed and tagged in
the code repository. It fixed the datasets, positive classes, method configuration,
baselines and search budgets, outer protocol, hypotheses H1--H4, the statistical
analysis plan, and the SHA-256 hashes of all source, script, test and requirement
files. Two deviations occurred, both of them engineering-only and documented with
their reasons and new hashes. D1: a parser fallback was needed because the default
ARFF parser fails under the installed pandas version for integer-like class labels.
D2: a unit test that asserted the held-out files were absent was conditioned on the
lock. Neither affects the method, the data or the metrics.

\subsection{Models}
Eleven models were compared. Logistic regression (L2) and a depth-limited decision
tree were tuned by grid search. XGBoost~\cite{chen2016}, LightGBM~\cite{ke2017},
CatBoost~\cite{prokhorenkova2018} and an MLP (scikit-learn~\cite{pedregosa2011}) were
tuned by a 20-iteration random search. XGBoost was also combined with post-hoc
Platt and isotonic calibration~\cite{zadrozny2002,niculescu2005}. All tuning used inner
stratified 3-fold cross-validation with log-loss scoring, so tuning itself
favoured calibration as much as discrimination. Two single-objective genetic-fuzzy
baselines share the encoding, operators and budget of the proposed model:
SO-GA-Fuzzy (classic), with evolved consequents and training accuracy as fitness,
represents prior practice; SO-GA-Fuzzy (hybrid), with the hybrid consequents of
\eqref{eq:consequents} and inner AUROC as fitness, isolates the effect of the
multi-objective search. The proposed MO-Fuzzy (NSGA-II) is described in
Section~\ref{sec:method}.

\subsection{Leakage-free evaluation}
Each dataset was evaluated by stratified 5-fold cross-validation repeated with
\resNSeeds{} seeds, giving \resFoldsPerDataset{} outer test folds per dataset and
model. Every fitted operation (imputation, scaling, one-hot encoding, tuning,
evolutionary search, $\lambda$ selection and operating-point selection) runs inside the
outer training fold, which addresses selection bias~\cite{cawley2010} and
leakage~\cite{kaufman2012}. Unit tests check that splits are disjoint and exhaustive;
that every estimator is fitted only on training rows (a spy records the rows passed to
\texttt{fit}); that fitted preprocessing does not change when test rows are
perturbed; that predictions are row-independent; that probabilities are valid; that
results are deterministic under seeding; and that Pareto non-dominance and hypervolume
are computed correctly.

\subsection{Metrics and statistics}
Discrimination is measured by AUROC and PR-AUC~\cite{saito2015}. Calibration is
measured by the Brier score~\cite{brier1950}, log-loss, ECE with 10 equal-width
bins~\cite{naeini2015,guo2017}, adaptive ECE~\cite{nixon2019}, and the calibration slope
and intercept of a logistic recalibration model~\cite{steyerberg2010,vancalster2016,huang2020}.
Interpretability is measured by distinct rules, mean antecedents per rule and total
literals; for the decision tree, leaves and the sum of leaf depths. Pareto quality is
measured by the normalised hypervolume~\cite{zitzler1999}, with objective scales
$(0.5, 0.25, R_{\max}L_{\max})$ and reference point $(1,1,1)$. We also record
single-threaded fit time and peak resident memory.
Following the preregistered plan, across-dataset comparisons use the Friedman
test~\cite{friedman1937} with Nemenyi critical differences~\cite{demsar2006}
(q$_{0.05}$ = \resCDq{} for $k=11$), and two-sided Wilcoxon signed-rank tests of
MO-Fuzzy against each model with Holm correction~\cite{holm1979}. With $N=6$ datasets the
smallest attainable two-sided Wilcoxon $p$ is 0.031, so these pairwise tests cannot
reach significance after correction. This limitation was declared in advance.
Per-dataset comparisons use the Nadeau--Bengio corrected resampled
$t$-test~\cite{nadeau2003} over the 25 paired folds, Holm-corrected within dataset and
metric. Reporting follows TRIPOD~\cite{collins2015} where applicable.

\subsection{Computation}
All models ran single-threaded (BLAS, OpenMP and model thread pools set to 1) on a
Windows 11 laptop with 16\,GB RAM, through a checkpointed runner that saves every fit
atomically and resumes after interruption. The full study (\resNFits{} fits) took
\resComputeHours{}\,h of fit time, and peak process memory never exceeded
\resMaxRSS{}\,MB.
